% -------------------------------------------------------------------------
% ÜBERHOLT (31.08.2026). Maßgeblich ist die eingereichte main.tex. Zahlen,
% Zeilennummern und Angaben zum Overleaf-Stand in dieser Datei stammen aus
% der Zeit VOR der Crosswalk-Korrektur vom 29.08.2026 (35 statt 36 Stadt-
% teile, vollständiger Neulauf) und wurden bewusst NICHT nachgezogen: die
% Datei ist als datierter Entwurf Teil der Arbeitsdokumentation.
% -------------------------------------------------------------------------
% =========================================================================
% KAPITEL 5.3 -- GLEITOBJEKTE
% 24.08.2026. Drei Objekte: zwei Quelltexte und die Merkmalstabelle.
%
% PLATZIERUNG
%   Quelltext 2 (lst:quoten)     -> nach Absatz 3 (Quoten)
%   Quelltext 3 (lst:krimindex)  -> nach Absatz 4 (Kriminalitaetsindex)
%   Tabelle (tab:codebook)       -> nach Absatz 8 (Zielgroessen).
%                                   4.1 verweist bereits darauf.
%
% ZUR TABELLE -- WAS GEGENUEBER DER BESTEHENDEN FASSUNG GEAENDERT IST
%   1) "flaechengewichtet" -> "bevoelkerungsgewichtet". acs_je_neighborhood
%      gewichtet mit total_population, nicht mit Flaeche. Betraf
%      Medianeinkommen und Medianmiete.
%   2) Wohnbevoelkerung war ebenfalls als "flaechengewichtet" gefuehrt --
%      sie wird ueber die Tracts SUMMIERT, nicht gewichtet gemittelt.
%   3) Spalte Skala: "absolut" ist kein Skalenniveau. Zaehlgroessen sind
%      verhaeltnisskaliert; die Spalte fuehrt jetzt durchgehend Verhaeltnis,
%      Intervall oder nominal. Schroeters Auflage vom 10.08. verlangt
%      ausdruecklich das Skalenniveau.
%   4) Risikogewerbe: "Handel, Gastgewerbe und Produktion" trifft die
%      Kategorien nicht. HIGH_RISK_COMMERCIAL = {RETAIL/ENT, PDR}, also
%      Handel und Unterhaltung sowie Produktion, Logistik und Reparatur.
%   5) Kriminalitaetsindex: Begriff "Standortquotient" beibehalten und der
%      Nenner benannt -- die Bevoelkerung, nicht die Gesamtzahl der Delikte.
%      Der Unterschied traegt den Absatz im Fliesstext.
%   Wertebereiche und Einheiten sind UNVERAENDERT uebernommen.
%
% KEINE deskriptive Statistik je Merkmal -- Auflage Schroeter 10.08.2026.
% Mittelwert, Median und Streuung stehen in Tabelle~\ref{tab:kennzahlen}
% in Abschnitt 4.2 und werden hier nicht wiederholt.
%
% Die Tabelle braucht tabularx und booktabs, beides ist in Kapitel 4 schon
% im Einsatz.
% =========================================================================


% ---------------------------- AB HIER KOPIEREN ---------------------------

\begin{lstlisting}[style=pythonstil,
  label={lst:quoten},
  caption={Anteilswerte aus Zählvariablen
           (\texttt{prep/s1\_daten.py}, Z.~559--584, gekürzt)}]
QUOTEN = [
    ("poverty_rate",     "poverty_below",         "poverty_universe_total"),
    ("bachelor_rate",    "bachelor_degree_count", "education_universe_total"),
    ("vacancy_rate",     "vacant_housing_units",  "total_housing_units"),
    ("pct_pre1940",      "pre1940_count",         "yrbuilt_count"),
    ("pct_residential",  "residential_count",     "parcel_count"),
]

def berechne_quoten(df: pd.DataFrame) -> pd.DataFrame:
    for name, zaehler, nenner in QUOTEN:
        z = pd.to_numeric(df[zaehler], errors="coerce").astype(float)
        n = pd.to_numeric(df[nenner], errors="coerce").astype(float)
        df[name] = (z / n.where(n > 0, np.nan)).round(4)
    return df
\end{lstlisting}


\begin{lstlisting}[style=pythonstil,
  label={lst:krimindex},
  caption={Relativer Kriminalitätsindex je Stadtteil und Monat
           (\texttt{prep/s1\_daten.py}, Z.~465--488, gekürzt)}]
    # Rollierendes Fenster, endend im VORMONAT: shift(1) VOR rolling()
    raster["delikte_fenster"] = (
        raster.groupby("neighborhood")["delikte"]
              .transform(lambda s: s.shift(1).rolling(CRIME_FENSTER_MONATE).sum()))
    raster = raster.dropna(subset=["delikte_fenster"])

    stadt = (raster.groupby(["jahr", "monat"])
                   .agg(stadt_delikte=("delikte_fenster", "sum"),
                        stadt_bev=("total_population", "sum")).reset_index())
    raster = raster.merge(stadt, on=["jahr", "monat"], how="left")

    raster["crime_rate_raw"] = (raster["delikte_fenster"]
                                / raster["total_population"] * 1000).round(3)
    raster["crime_index"] = (raster["crime_rate_raw"]
                             / (raster["stadt_delikte"] / raster["stadt_bev"] * 1000)
                             ).round(4)
\end{lstlisting}


\begin{table}[htbp]
\centering
\caption{Merkmale und Zielgrößen des Analysedatensatzes}
\label{tab:codebook}
\vspace{4pt}
\small
\begin{tabularx}{\textwidth}{@{}l l l X@{}}
\toprule
\textbf{Größe} & \textbf{Skala} & \textbf{Wertebereich} & \textbf{Quelle und Bildung} \\
\midrule
\addlinespace
\multicolumn{4}{@{}l}{\textbf{Merkmale --- sozioökonomisch}} \\
\addlinespace
Medianeinkommen & Verhältnis & 20.562 bis 246.635 USD je Jahr & ACS B19013, bevölkerungsgewichtetes Mittel der Tract-Werte \\
Armutsquote & Verhältnis & 0,007 bis 0,361 & ACS B17001, Anteil an der Bezugsbevölkerung \\
Akademikerquote & Verhältnis & 0,147 bis 0,553 & ACS B15003, Anteil an der Bezugsbevölkerung \\
Medianmiete & Verhältnis & 722 bis 3.501 USD je Monat & ACS B25064, bevölkerungsgewichtetes Mittel der Tract-Werte \\
Leerstandsquote & Verhältnis & 0,005 bis 0,237 & ACS B25002, Anteil an allen Wohneinheiten \\
\addlinespace
\multicolumn{4}{@{}l}{\textbf{Merkmale --- kriminalitätsbezogen}} \\
\addlinespace
Kriminalitätsindex, logarithmiert & Intervall & $-2{,}98$ bis $2{,}61$ & SFPD, Standortquotient der Deliktrate je Einwohner über die zwölf Vormonate; 0 ist der Stadtdurchschnitt \\
\addlinespace
\multicolumn{4}{@{}l}{\textbf{Merkmale --- baulich}} \\
\addlinespace
Altbauanteil vor 1940 & Verhältnis & 0,128 bis 0,908 & Land Use 2020, Anteil an den Parzellen mit bekanntem Baujahr \\
Wohngebäudeanteil & Verhältnis & 0,133 bis 0,971 & Land Use 2020, Anteil an allen Parzellen \\
Risikogewerbeanteil & Verhältnis & 0,000 bis 0,222 & Land Use 2020, Flächenanteil der Nutzungsarten Handel und Unterhaltung sowie Produktion, Logistik und Reparatur \\
\addlinespace
\multicolumn{4}{@{}l}{\textbf{Merkmale --- Größenkontrolle und Saison}} \\
\addlinespace
Wohnbevölkerung, logarithmiert & Intervall & 7,77 bis 11,30 & ACS B01003, über die Tracts summiert und logarithmiert \\
Saison, Sinus & Intervall & $-1{,}000$ bis $1{,}000$ & abgeleitet, $\sin(2\pi \cdot \text{Monat}/12)$ \\
Saison, Kosinus & Intervall & $-1{,}000$ bis $1{,}000$ & abgeleitet, $\cos(2\pi \cdot \text{Monat}/12)$ \\
\addlinespace
\multicolumn{4}{@{}l}{\textbf{Zielgrößen und Expositionsgröße}} \\
\addlinespace
Einsatzhäufigkeit & Verhältnis & 0 bis 451 Einsätze je Monat & SFFD, Zählung je Stadtteil und Monat \\
Einsätze je 1.000 Einwohner & Verhältnis & 0,00 bis 67,66 & abgeleitet, auf die Wohnbevölkerung normiert \\
überwiegende Einsatzart & nominal & Brand, Rettung, technische Hilfe, Fehlalarm & SFFD, größter der vier Anteile desselben Monats \\
Wohnbevölkerung & Verhältnis & 2.357 bis 81.209 Personen & ACS B01003, über die Tracts summiert; Exposition \\
\bottomrule
\end{tabularx}
\end{table}

% ---------------------------- BIS HIER KOPIEREN --------------------------
